\section{Introduction}
Tropical cyclones can alter forest structure and bird habitat over time scales extending well beyond the storm itself. Research in Puerto Rico has documented both forest damage and subsequent bird responses, providing a basis for asking how disturbance-related change varies across an island landscape~\cite{uriarte2019,wunderle1995}. Recent work on the Puerto Rico Plain Pigeon documents a prolonged post-hurricane population bottleneck after accounting for detection, while remote sensing of Puerto Rican mangroves shows that canopy damage and vegetation-index recovery vary among sites~\cite{rivera2025,howe2025}. These studies motivate environmental comparisons without assuming that vegetation recovery and bird recovery follow the same trajectory. For an endemic resident species, a rise in records after a period of low reporting needs to be examined in its environmental context. Analysis must locate the increase, test whether elevation and vegetation greenness are associated with previously occupied sites remaining occupied in the next period, and assess whether the relationships can guide monitoring elsewhere. Addressing these questions requires keeping the observed pattern, its ecological interpretation and its predictive usefulness distinct.

The Puerto Rican Emerald (\Rm) offers an opportunity to examine this sequence. Ortiz-Andrade and Rojas-Perez~\cite{ortiz2026} combined eBird checklists, elevation and satellite vegetation greenness to describe three periods defined relative to Hurricanes Mar\'ia and Fiona. Their checklist-level models linked reporting to environmental gradients, while a dynamic occupancy model addressed occupancy persistence under imperfect detection. Their public reproducibility package permits a reanalysis of these complementary lines of evidence. It also permits a distinction between reproducing a reported number and checking whether the implementation represents the intended statistical model.

The first interpretive challenge concerns the response being measured. An eBird checklist describes a reporting event, whose probability depends on species presence, survey effort and other observation conditions~\cite{sullivan2009,johnston2021}. Backstrom et al.~\cite{backstrom2025} further show that spatial and temporal sampling biases can interact in eBird data and alter trend estimates, making the environmental composition of the sample relevant to period comparisons. A change in the proportion of positive checklists is therefore not, by itself, a change in population size. Elevation-specific predictions can reveal differences hidden by island-wide averages, but they remain conditional reporting probabilities. Dynamic occupancy models estimate latent site occupancy and transitions between occupied and unoccupied states, conditional on repeat-survey and detection assumptions~\cite{mackenzie2002,mackenzie2003,guillera2017}. A recent Prairie Warbler study used long-term point counts and dynamic occupancy models to distinguish colonization and local extinction responses to habitat management~\cite{gordon2024}, illustrating why site retention and entry into occupancy should be evaluated separately. These quantities must be interpreted on their own terms. In particular, a high occurrence score or its complement does not measure a local extinction transition.

The second challenge concerns the reliability of the observation process used to infer those states. Environmental associations with persistence are informative only insofar as observations and covariates are correctly aligned. A reproducible script may reproduce an unintended arrangement just as consistently as an intended one. Examining implementation errors helps assess the reliability of ecological interpretations. A correction may retain the environmental association while changing the evidence for a detection difference, making separate evaluation of these explanations necessary.

The third challenge is geographic prediction. A relationship that summarizes the sampled locations may not retain the same value in held-out geographic blocks. Repeated visits and geographically clustered checklists make the validation design relevant to the intended monitoring application~\cite{roberts2017,valavi2019,ploton2020}. Random and spatial partitions ask different questions, and good performance under one does not establish performance under the other. Recent comparisons of species-distribution models confirm that random validation can overestimate transfer performance~\cite{spacetime2025}; tests of spatial versus temporal environmental relationships also show that transferability can vary within a species' range~\cite{lovell2025}. Moreover, predicting checklist reports is not an independent validation of a latent occupancy model. Ecological explanation and reporting prediction contribute different evidence to a monitoring decision~\cite{elith2009}. In bird forecasting, Ovaskainen et al.~\cite{ovaskainen2026} combined citizen-science recordings with standardized intervals and permanent point-count networks, illustrating how prediction can be supported by a repeatable observation design.

We examine reporting rebound across elevation and greenness conditions, test whether these environments also have higher occupancy persistence probabilities, and use the results to assess how environmental differences should be monitored. First, missingness composition and standardized reporting predictions separate environmental representation from within-stratum period change. Second, corrected period-label alignment in the dynamic occupancy model allows reporting trajectories to be compared with occupancy persistence probability, colonization and repeat reports in corresponding environmental cells. Finally, spatial and temporal validation and equal-budget stratified ranking assess whether reporting predictions can supply clues for revisits in target environments. Conventional benchmarks and checks of the exploratory neural workflow serve this monitoring question. As a secondary analysis of the same public observations, the study connects the environmental distribution of rebound, corrected persistence associations and practical tradeoffs in monitoring coverage.
